\documentclass[11pt,a4paper]{article}
\usepackage[utf8]{inputenc}
\usepackage[T1]{fontenc}
\usepackage{amsmath, amssymb}
\usepackage{graphicx}
\usepackage{hyperref}
\begin{document}
% [NODE_ID: macro_1]
ELSEVIER

Listas de contenido disponibles en ScienceDirect

Revista Trimestral de Economía y Finanzas

página de inicio de la revista: www.elsevier.com/locate/qref

\% [Imagen omitida]

Desalineaciones del tipo de cambio real y crecimiento económico en Túnez: Nueva evidencia de un análisis de umbral de ajustes asimétricos\*

\% [Imagen omitida]

Thouraya Hadj Amor <sup>a,e</sup>, Ridha Nouira <sup>b</sup>, Christophe Rault <sup>c,\*,1</sup>, Anamaria Diana Sova <sup>c</sup>

- a Departamento de Administración de Empresas, Universidad de Shaara, Arabia Saudita
- b ISFF y LAMIDED Universidad de Sousse, Túnez
- c LEO, Université d'Orléans, Rue de Blois-B.P. 6739, 45067 Orléans Cedex 2, Francia
- <sup>d</sup> Brunel University London, Reino Unido
- e Departamento de Economía- EAS. Universidad de Monastir, Túnez

INFORMACIÓN DEL ARTÍCULO

Historial del artículo:
Recibido el 18 de febrero de 2021
Recibido en forma revisada el 22 de enero de 202:
Aceptado el 26 de enero de 2023
Disponible en línea el 27 de enero de 2023

Clasificación IFI

F31

F43

C10

Palabras clave: Crecimiento Desalineación del tipo de cambio No linealidades NARDL Umbrales

% [NODE_ID: macro_2]
RESUMEN

En este artículo, nos centramos tanto en el impacto de los desalineamientos del tipo de cambio real (TCR) en el crecimiento económico como en la evaluación de los umbrales de los desalineamientos en el caso de Túnez. Nuestra contribución a la literatura es triple. Primero, reevaluamos los principales factores que impulsan el TCR de equilibrio utilizando un enfoque de modelo dinámico. Segundo, investigamos los posibles efectos asimétricos del desalineamiento real en el crecimiento económico de Túnez utilizando el estimador de rezagos distribuidos autorregresivos no lineales (NARDL) desarrollado por Shin et al. (2014), que tiene en cuenta la no linealidad en la relación entre los desalineamientos del tipo de cambio y el crecimiento económico. Tercero, utilizando un modelo de umbral inspirado en Hansen (2000), identificamos el nivel de desalineamientos del dinar tunecino que puede mejorar (o dificultar) el crecimiento económico. Nuestros principales hallazgos son los siguientes: i) el dinar tunecino ha experimentado períodos de sobrevaloración y subvaloración durante el período 2001-2016; ii) existe un impacto negativo de la sobrevaloración real del dinar en el rendimiento del crecimiento de Túnez, mientras que una subvaloración no tiene un impacto significativo; iii) existe un umbral para los efectos del desalineamiento del TCR en el crecimiento. Específicamente, la subvaloración real promueve el crecimiento hasta que se alcanza un umbral estimado del 10.02\% de desviación del TCR de equilibrio, mientras que una sobrevaloración puede impulsar el crecimiento hasta que se alcanza un umbral del 9.02\%. Estos hallazgos tienen importantes implicaciones desde la perspectiva de un formulador de políticas. © 2023 Junta de Fideicomisarios de la Universidad de Illinois, Publicado por Elsevier Inc. Todos los derechos reservados.

% [NODE_ID: macro_3]
1. Introducción

Después de 2010, el proceso de convergencia económica tunecina experimentó una desaceleración notable. De hecho, este período estuvo marcado por una inestabilidad del clima social y un deterioro considerable de la seguridad. El aumento de la demanda, que ejerció presión sobre los precios, se explicó principalmente por el aumento del consumo privado apoyado por la creación exagerada de empleos en el sector público y

Direcciones de correo electrónico: t.amor@su.edu.sa, hathouraya@yahoo.fr (T.H. Amor)
nouira.ridha75@gmail.com (R. Nouira), chrault@hotmail.com (C. Rault),
AnamariaDiana Soya@brunel ac uk (A.D. Soya)

Sitio web: http://chrault3.free.fr

así el aumento de la masa salarial. Por el contrario, el sector privado se ha caracterizado por un deterioro de su actividad económica.

Los disturbios sociales, así como el deterioro de la seguridad, también intensificados por la crisis libia, afectaron los pagos externos. En este contexto, el déficit por cuenta corriente aumentó considerablemente, pasando en promedio del 3.1\% del PIB entre 2006 y 2010 al 9.1\% del PIB entre 2013 y 2017. El déficit comercial se deterioró del 13.2\% del PIB en 2010 al 16\% del PIB en 2017. Esto se debió principalmente a problemas de seguridad en 2011, la desaceleración de la actividad productiva de las industrias extractivas, en particular el sector industrial del fosfato en Gafsa, y las importaciones incontroladas de bienes de consumo. Además, la disminución de las inversiones en exploración de petróleo y gas ha aumentado el déficit de la balanza energética en los últimos años.

Túnez entró en la crisis de la pandemia de COVID-19 con una lenta tasa de crecimiento económico y niveles de deuda crecientes, y se ha visto fuertemente afectado por la crisis debido a su gran dependencia de Europa. La profundidad del impacto de la pandemia en la economía tunecina se hizo

\textsuperscript{*} Agradecemos a dos revisores anónimos por sus muy útiles sugerencias en versiones anteriores del artículo.

\textsuperscript{*} Autor de correspondencia.

más evidente a finales de 2020. De hecho, el PIB real se contrajo un 8.8\% en 2020 debido a la disminución general de la actividad económica, caracterizada por una disminución de la producción en los principales sectores, en el sector de servicios (incluido el turismo), la inversión y las exportaciones. El desempleo aumentó del 15\% antes de la pandemia al 17.8\% a finales del primer trimestre de 2021. La cuenta corriente se estabilizó en un déficit del 8.1\% del PIB en 2020. Las exportaciones totales en 2020 disminuyeron un 11.7\% en comparación con 2019 (Banco Mundial, 2019, 2021).

En este contexto, el equilibrio macroeconómico y la competitividad externa de cualquier país siguen dependiendo de un ajuste del valor real del tipo de cambio y su valor de equilibrio, una necesidad que se explica por las consecuencias perjudiciales de los desajustes recurrentes y persistentes, que pueden conducir a desequilibrios estructurales en el mercado laboral y el mercado de bienes transables y no transables (Elfathaoui, 2019). De hecho, la evolución del proceso de integración financiera internacional ha generado una creciente desconexión entre la variabilidad del RER y el crecimiento (Hadj Amor, 2009; Béreau et al., 2012). Un gran volumen de literatura ha destacado la importancia de observar de cerca el tipo de cambio para países en transición como Túnez, y que también puede utilizarse como medida reguladora para mitigar, si no eliminar, los efectos adversos de un tipo de cambio sobrevalorado (Derbali, 2021).

En el contexto tunecino, la cuestión del desajuste del tipo de cambio real, definido como la desviación del tipo de cambio real efectivo de su nivel de equilibrio, y su impacto en el crecimiento económico ha sido objeto de debate durante algún tiempo por muchas razones, un debate que se ha acentuado después de la adopción de un régimen cambiario más flexible. De hecho, desde 1990, el Banco Central de Túnez (BCT) ha optado por una política de tipo de cambio flexible acompañada de una mayor liberalización económica en Túnez, con el objetivo de dar a la política cambiaria un papel más activo. El objetivo era mejorar la competitividad de la economía nacional al tiempo que se estabilizaba el dinar tunecino (Balazs et al., 2014).

% [NODE_ID: macro_4]
El último deterioro de las condiciones económicas tunecinas, como se señaló anteriormente, puede apuntar hacia una posible distorsión del RER. Dada la gravedad de las consecuencias de un posible desajuste de la economía, las autoridades monetarias tunecinas deben priorizar y analizar la cuestión del desajuste del RER tunecino y sus efectos en el crecimiento económico.

Túnez, como pequeña economía abierta, se ha comprometido desde 2011 con un Nuevo Marco Operativo, basado en la autonomía monetaria y la transición de un tipo de cambio de paridad móvil (crawling peg) a una banda de fluctuación móvil (crawling band). Además, en los últimos años, el dinar, la moneda nacional de Túnez, ha registrado una fuerte depreciación de su valor frente a otras monedas extranjeras. Más recientemente, especialmente después de la Revolución de Jazmín de 2011, el dinar tunecino rompió su récord histórico y se ha depreciado frente al euro y el dólar estadounidense. Si bien la depreciación del dinar, en un modelo de equilibrio parcial de la balanza comercial, puede aumentar las exportaciones netas con algunos países socios, bajo el supuesto de que hay espacio para que el sector exportador crezca (Bahmani-Oskooee et al., 2017, 2019), sus beneficios podrían no superar el impacto negativo de una subvaluación en la economía real, dado que el valor de la moneda es inferior al que debería ser en equilibrio o está severamente depreciado.

Desde octubre de 2017, y para luchar contra la disminución de los activos en moneda extranjera, los agentes económicos que importan bienes de consumo no esenciales no tienen derecho a utilizar préstamos bancarios para estas operaciones. La deuda externa de Túnez aumentó considerablemente hasta alcanzar el 70\% del PIB a finales de 2016. Se esperaba que dicha deuda externa soportara varios choques, excepto la significativa depreciación del tipo de cambio real (FMI, 2017). En noviembre de 2017, las reservas de divisas estaban al nivel de cubrir 90 días de importaciones de bienes y servicios. La deuda pública tunecina, de la cual el 70\% es externa, ha aumentado del 72\% del PIB en 2019 al 87\% del PIB en 2020, lo que está muy por encima del punto de referencia de carga de deuda de los mercados emergentes del 70\% del PIB (Banco Central de Túnez (BCT) e Instituto Nacional de Estadística (INE), 2020).

Aunque algunos estudios han demostrado los costos de la subvaluación del tipo de cambio en el crecimiento económico (Grekou, 2015, 2018a; Chen, 2017; Morvillier, 2020), muchos países en desarrollo han dependido de un tipo de cambio real subvaluado para impulsar sus exportaciones (Levy-Yeyati \& Sturzzenegger, 2007; Elbadawi et al., 2012; Rodrik, 2009; Korinek \& Serven, 2010; Iyke \& Odhiambo, 2015; Habib et al., 2017; Mazorodze \& Tewari, 2018), pero a raíz de la crisis financiera global ha surgido un intenso debate sobre los méritos asociados de estas estrategias de crecimiento impulsadas por las exportaciones. Dado que las perspectivas económicas globales son más débiles que en el pasado, existe una mayor incertidumbre sobre la capacidad de las economías avanzadas para seguir absorbiendo las exportaciones de los países en desarrollo. Además, si el tipo de cambio real se mantiene demasiado bajo durante demasiado tiempo, es probable que las estrategias de crecimiento impulsadas por las exportaciones basadas en la subvaluación de la moneda generen costos para la economía agregada. En este contexto, la relación entre el desajuste del tipo de cambio real (REM), es decir, la desviación entre el Tipo de Cambio Real (RER) y su nivel de equilibrio, y la economía real, ha sido el foco de las recientes discusiones de política macroeconómica en Túnez. Ha habido muchos debates entre los analistas económicos tunecinos sobre las causas de la depreciación del RER y en qué medida refleja cambios en los fundamentos macroeconómicos de Túnez, lo que subraya la necesidad de reevaluar los principales determinantes del RER de equilibrio, el grado de distorsión del dinar tunecino y en qué medida el desajuste de la moneda tunecina afecta su desempeño de crecimiento.

Teóricamente, existe un debate en curso sobre si los desajustes del RER impiden o facilitan el crecimiento económico. La literatura relacionada gira en torno a dos ejes principales: a través del Consenso de Washington (WC), Williamson (1990) defiende un RER cercano a su nivel de equilibrio y consistente a medio plazo con objetivos macroeconómicos que estimulan el crecimiento. Por el contrario, la hipótesis del crecimiento impulsado por las exportaciones destaca la naturaleza asimétrica de los desajustes. De hecho, el crecimiento económico se ve obstaculizado por las sobrevaluaciones, mientras que se ve impulsado por la subvaluación de la moneda.

% [NODE_ID: macro_5]
La subvaloración de la moneda promueve el crecimiento a través de una acumulación de capital más rápida, cambio tecnológico y derrames de información a otras empresas e industrias en la economía. De hecho, la literatura empírica reciente que incorpora no linealidades en el análisis del nexo RER-crecimiento sugiere que las sobrevaloraciones y subvaloraciones pueden afectar de manera diferente al crecimiento (Hadj Amor \& Sarkar, 2009; Hadj Amor, 2009; Nouira \& Sekkat, 2012; Schroder, 2013), mientras que el nivel del desalineamiento del RER también importa (Aguirre \& Calderon, 2005; Couharde \& Sallenave, 2013; Tipoy et al., 2018).

Algunos estudios centrados en Túnez han discutido las cuestiones del RER de equilibrio y el desalineamiento. Charfi (2008) estimó el tipo de cambio real de equilibrio del Dinar de 1983 a 2000, utilizando datos trimestrales, basándose en las siguientes variables fundamentales: términos de intercambio; entradas netas de capital; y el diferencial de productividad. Los resultados mostraron que el dinar tunecino estaba sobrevalorado antes de la devaluación de 1986 y se acercó a su valor de equilibrio durante los años 90. A principios de este siglo (2000), las autoridades permitieron una mayor fluctuación del tipo de cambio real efectivo. Amaira (2021) estudió el desalineamiento del tipo de cambio real efectivo (REER) durante el período 1986-2015. Los resultados mostraron que el nivel de equilibrio del tipo de cambio a largo plazo depende de la productividad, los términos de intercambio y el gasto gubernamental, y que el desalineamiento es muy bajo, especialmente desde la implementación del plan de ajuste estructural. Derbali (2021) utilizó el enfoque del tipo de cambio de equilibrio conductual (BEER) para estimar el tipo de cambio de equilibrio tunecino utilizando modelos autorregresivos vectoriales y modelos de corrección de errores vectoriales para datos trimestrales durante el período 1990-2020. Los resultados mostraron una baja sensibilidad del RER a los choques monetarios y comerciales y una convergencia de la serie del tipo de cambio real de su trayectoria a su valor objetivo a largo plazo. Benzid (2021) evaluó los desalineamientos en el tipo de cambio real efectivo de Túnez con énfasis en su persistencia a lo largo del

\% [Imagen omitida]

Fig. 1. Desalineamiento del RE de Túnez.

Fuente: Cálculo del autor (RER en azul y ERER en verde, eje izquierdo: Mis en rojo. eje derecho)

Apéndice. Resulta que, excepto por los términos de intercambio (*ToT*), todas las demás variables del modelo son no estacionarias en niveles. Además, se rechazan las hipótesis nulas de una raíz unitaria para todas las series en primeras diferencias. Esto implica que todas las variables son I(1) excepto los términos de intercambio, que es estacionaria en nivel, I(0). Para verificar si los resultados de la prueba ADF-GLS son robustos, realizamos la prueba de raíz unitaria con quiebres estructurales de Zivot y Andrews (1992). Los resultados reportados en la Tabla A.4 del Apéndice confirman los resultados de la prueba ADF-GLS. Pesaran y Smith (1995), y Pesaran y Shin (1999) muestran que el enfoque ARDL puede usarse para el análisis a largo plazo, y que la metodología ARDL es válida incluso si las variables son I(0) o I(1). Por lo tanto, aplicamos un modelo autorregresivo de rezagos distribuidos (ARDL) o enfoque de prueba de límites (Pesaran et al., 2001) para estimar el tipo de cambio real de equilibrio en Túnez. Según Pesaran et al. (2001), este enfoque es más adecuado cuando el número de observaciones temporales es pequeño. Antes de realizar la estimación del modelo ARDL, es esencial escribir la Ec. (1) como un modelo ARDL condicional, de la siguiente manera:

$$\begin{split} \Delta IRER_t &= \mu + \sum_{i=1}^p \delta_{1i} \, \Delta IRER_{t-i} + \sum_{i=0}^p \delta_{2i} \Delta IToT_{t-i} + \sum_{i=0}^p \delta_{3i} \Delta IGov_{t-i} \\ &+ \sum_{i=0}^p \delta_{4i} \Delta IDef_{t-i} (eestimteeqelong - runandshort - run \\ &+ \sum_{i=0}^p \delta_{5i} \Delta IU_{t-i} + \sum_{i=0}^p \delta_{6i} \Delta IBalSam_{t-i} + \gamma_1 IIRER_{t-1} \\ &+ \gamma_2 IToT_{t-1} + \gamma_3 IGov_{t-1} + \gamma_4 IDef_{t-1} + \gamma_5 IU_{t-1} + \gamma_6 IBalSam_{t-1} + \varepsilon_t \end{split}$$

Todas las variables se expresan en logaritmo, △ es el operador de primera diferencia. La relación a largo plazo entre el RER y sus determinantes puede identificarse basándose en la prueba de Wald (estadístico F). Este estadístico determina la significancia conjunta de las variables de nivel rezagadas bajo la hipótesis nula. Si los coeficientes no son conjuntamente todos iguales a cero, entonces implica la existencia de cointegración, es decir, una relación a largo plazo entre el RER y sus determinantes. Una vez una cointegración

% [NODE_ID: macro_6]
establecida la relación, procedemos a la estimación de los parámetros a largo y corto plazo del modelo. Específicamente, estimamos los coeficientes a largo plazo utilizando la Ec. (3), como se muestra a continuación:

$$\begin{split} \mathit{IRER}_t &= \mu + \gamma_1 \mathit{IRER}_{t-1} + \gamma_2 \mathit{IToT}_{t-1} + \gamma_3 \mathit{IGov}_{t-1} + \gamma_4 \mathit{IDef}_{t-1} \\ &+ \gamma_5 \mathit{IBalSam}_{t-1} + \varepsilon_t \end{split} \tag{3}$$

Los coeficientes dinámicos a corto plazo pueden obtenerse estimando el modelo de corrección de error vectorial (VEC) correspondiente, como se muestra a continuación – Ec. (4):

$$\begin{split} \Delta IRER_t &= \mu + \sum_{i=1}^p \delta_{1i} \Delta IRER_{t-i} + \sum_{i=0}^p \delta_{2i} \Delta IToT_{t-i} + \sum_{i=0}^p \delta_{3i} \Delta IGov_{t-i} \\ &+ \sum_{i=0}^p \delta_{4i} \Delta IDef_{t-i} (eestimteeqelong - runandshort - run \\ &+ \sum_{i=0}^p \delta_{5i} \Delta IBalSam_{t-i} + \eta ECM_{t-1} + \varepsilon_t \end{split}$$

donde $ECM_{t-1}$ es el término de corrección de error rezagado un período estimado de la Ec. (4). $\eta$ es el coeficiente para la velocidad de ajuste del modelo a la situación de equilibrio, que se espera que tenga un signo negativo.

Los límites ARDL estimados muestran que el estadístico F de la prueba de Wald, 4.82, es significativo al nivel del 1\%. De hecho, al comparar los estadísticos F con los valores críticos extraídos de Pesaran et al (2001), el estadístico F es superior al valor crítico del límite superior de 4.68 al nivel de significancia del 1\%. Esto implica que la hipótesis nula (Ho: $\gamma_1 = \gamma_2 = \gamma_3 = \gamma_4 = \gamma_5 = \gamma_6 = 0$ ) de no relación a largo plazo puede ser rechazada por los datos. Por lo tanto, hay evidencia de cointegración entre el RER y sus determinantes. A continuación, estimamos los coeficientes a largo plazo del modelo de tipo de cambio real de equilibrio. Para ello, estimamos modelos ARDL (Ec. 2). La especificación de nuestro modelo ARDL se determina mediante el criterio de información de Akaike (AIC). Los resultados de la Tabla 1 muestran que los coeficientes para las variables *Gov*, *Balassa*, *TOT* y *Def* son estadísticamente significativos y tienen los signos esperados. El coeficiente estimado para la velocidad de ajuste ( $\eta$ ) es negativo y es estadísticamente significativo al nivel del 5\%.

Los resultados empíricos muestran un impacto positivo y significativo de los términos de intercambio, el gasto gubernamental y el efecto Balassa–Samuelson, y un efecto negativo y significativo del déficit comercial. Para los términos de intercambio, una mejora del 10\% conduce a una apreciación del ERER del 2.4\%. Esto implica que el efecto ingreso supera al efecto sustitución. BalSam es una variable que busca

<sup>&</sup>lt;sup>4</sup> La prueba eficiente de raíz unitaria sugerida por Kwiatkovski-Phillips-Schmidt-Shir (1992), extendida por Carrion-i-Silvestre y Sanso (2006), también ha sido implementada, lo que lleva a resultados similares.

<sup>&</sup>lt;sup>5</sup> Nótese, sin embargo, que rechazar la nula ( $\gamma_1 = \gamma_2 = \gamma_3 = \gamma_4 = \gamma_5 = \gamma_6 = 0$ ) no implica necesariamente que todos los coeficientes $\gamma_i \neq 0$ , para i = 1, ..., 6. Por ejemplo, en la Tabla 1, el coeficiente de desempleo es igual a cero, lo que explica por qué esta variable no está incluida en la ecuación (3). Además, una investigación más profunda también revela que el desempleo no tiene influencia sobre el RER a corto plazo, véase la ecuación (4).

\end{document}